%% NuMI CC1pi internal note
%% Companion to BNBInternalNoteCC1pi_V0.91.pdf
%% Build:  pdflatex NuMIInternalNoteCC1pi.tex   (twice, for the ToC)

\documentclass[11pt,a4paper]{article}

% ---------------------------------------------------------------- packages --
\usepackage[T1]{fontenc}
\usepackage[utf8]{inputenc}
\usepackage{lmodern}
\usepackage[english]{babel}

\usepackage[
  a4paper,
  top=2.5cm, bottom=2.5cm, left=2.7cm, right=2.7cm,
  headheight=14pt
]{geometry}

\usepackage{amsmath}
\usepackage{amssymb}
\usepackage{graphicx}
\usepackage{booktabs}
\usepackage{tabularx}
\usepackage{array}
\usepackage{longtable}
\usepackage{caption}
\usepackage{xcolor}
\usepackage{enumitem}
\usepackage{fancyhdr}
\usepackage{microtype}
\usepackage[hidelinks,breaklinks]{hyperref}

% ------------------------------------------------------------------ styling --
\definecolor{warnbg}{RGB}{253,246,227}
\definecolor{warnrule}{RGB}{181,137,0}
\definecolor{codebg}{RGB}{245,245,245}

\setlength{\parskip}{0.5em}
\setlength{\parindent}{0pt}

% This note is dense with long unbreakable identifiers (file paths, config
% names, branch names) set in \texttt. Without these, they overrun the right
% margin. \sloppy trades a little interword spacing for correct margins, and
% \code below permits breaks at the separators those identifiers actually use.
\sloppy
\setlength{\emergencystretch}{3em}

\captionsetup{font=small,labelfont=bf}

\pagestyle{fancy}
\fancyhf{}
\lhead{\small NuMI CC1$\pi$ internal note}
\rhead{\small Draft v0.1}
\cfoot{\thepage}
\renewcommand{\headrulewidth}{0.4pt}

% Callout box. mdframed is not available in this TeX installation, so this is
% hand-rolled from a saved box inside a \colorbox. Using a minipage (rather
% than a tabular with \cellcolor, which would additionally require colortbl)
% keeps multi-paragraph content working inside the callout.
\newsavebox{\calloutbox}
\newenvironment{callout}[1]{%
  \par\medskip
  \begin{lrbox}{\calloutbox}%
  \begin{minipage}{\dimexpr\linewidth-2\fboxsep\relax}%
  \setlength{\parskip}{0.5em}%
  \textbf{#1}\par
}{%
  \end{minipage}%
  \end{lrbox}%
  \noindent\colorbox{warnbg}{\usebox{\calloutbox}}%
  \par\medskip
}

% \code: typewriter, and allowed to break across lines. hyphenat and seqsplit
% are not available in this TeX installation. Re-enabling \hyphenchar for the
% typewriter font lets TeX break long identifiers; combined with \sloppy above
% this keeps them inside the margin. Note that a break inserts a hyphen, so
% for identifiers where an inserted hyphen would be ambiguous, prefer \path.
\newcommand{\code}[1]{{\small\ttfamily\hyphenchar\font=`\-\relax #1}}

% \path: typewriter with NO hyphenation, for file paths and anything where a
% spurious hyphen would change the meaning.
\newcommand{\pathname}[1]{\texttt{\small #1}}
\newcommand{\numu}{\ensuremath{\nu_\mu}}
\newcommand{\numubar}{\ensuremath{\bar{\nu}_\mu}}
\newcommand{\pmu}{\ensuremath{p_\mu}}
\newcommand{\ppi}{\ensuremath{p_\pi}}
\newcommand{\GeVc}{\ensuremath{\mathrm{GeV}/c}}

% -------------------------------------------------------------------- title --
\title{\vspace{-1.5cm}\textbf{Muon Neutrino Charged Current Single Pion\\
Cross Section on Argon using Run 1 NuMI Data}}
\author{MicroBooNE Internal Note --- Draft}
\date{Version 0.1}

\begin{document}
\maketitle
\thispagestyle{fancy}

\begin{callout}{DRAFT STATUS --- READ FIRST}
This note mirrors the structure of the BNB CC1$\pi$ internal note
(\code{BNBInternalNoteCC1pi\_V0.91.pdf}, J.~P.~Detje, 4 February 2025) for the
NuMI Run~1 FHC measurement.

Sections~\ref{sec:overview}--\ref{sec:extraction} describe the analysis
\emph{as it is actually configured in the code today}, and every number in
them was read from the source or the configuration files.
Sections~\ref{sec:fakedata}--\ref{sec:results} are \emph{structural
placeholders}. No cross-section values, efficiencies, purities, $\chi^2$
values or figures appear anywhere in this note, because the analysis has not
been re-run since the corrections described in
Section~\ref{sec:corrections}. Nothing here should be quoted as a result.

Three items must be settled before this note can carry results at all; see
Section~\ref{sec:open}.
\end{callout}

\begin{abstract}
\noindent
This note describes the extraction of the \numu\ charged-current single
charged pion cross section on argon using MicroBooNE Run~1 FHC data from the
NuMI beam. It is the NuMI counterpart of the BNB analysis of the same final
state and is intended to be read alongside that note, which contains fuller
descriptions of the selection variables and the boosted decision trees. The
signal definition, binning, selection, uncertainty treatment and unfolding
procedure are documented; the fake-data tests and results sections are
reserved pending a complete re-run of the analysis chain.
\end{abstract}

\tableofcontents
\newpage

% =============================================================================
\section{Overview}
\label{sec:overview}

The measurement is performed with the common \code{xsec\_analyzer} framework:
event selection via the \code{CC1mu1piXp} selection class, universe
construction via \code{univmake}, and unfolding via the Wiener-SVD
implementation in \code{CrossSectionExtractor}.

Section~\ref{sec:software} describes the software and samples,
Section~\ref{sec:selection} the signal definition and selection,
Section~\ref{sec:unc} the uncertainty treatment,
Section~\ref{sec:extraction} the cross-section extraction and unfolding,
Section~\ref{sec:fakedata} the fake-data tests and
Section~\ref{sec:results} the results.

The principal differences from the BNB analysis are summarised in
Table~\ref{tab:diffs}.

\begin{table}[h]
\centering
\small
\begin{tabularx}{\linewidth}{@{}lXX@{}}
\toprule
 & \textbf{BNB analysis} & \textbf{This analysis} \\
\midrule
Beam                & BNB & NuMI FHC \\
Runs                & 1--5 & 1 only \\
Flux systematic     & \code{weight\_flux\_all} & \code{weight\_ppfx\_all} (PPFX) \\
Beamline geometry   & not applicable & NuMI-specific; \textbf{not currently
                      included}, see Section~\ref{sec:gaps} \\
Proton multiplicity & --- & Xp (any number, no threshold) \\
Unfolding           & Wiener-SVD & Wiener-SVD, D'Agostini as cross-check \\
\bottomrule
\end{tabularx}
\caption{Principal differences from the BNB CC1$\pi$ analysis.}
\label{tab:diffs}
\end{table}

% =============================================================================
\section{Software and Samples}
\label{sec:software}

\subsection{Framework}

The analysis runs entirely within \code{xsec\_analyzer}. The chain is:

\begin{enumerate}[itemsep=0.2em,topsep=0.3em]
  \item \code{ProcessNTuples} --- applies the \code{CC1mu1piXp} selection to
        raw PeLEE ntuples and writes \code{stv\_tree} output
        (\code{run\_process\_ntuples.sh}).
  \item \code{univmake} --- builds systematic universe histograms from a bin
        configuration (\code{run\_universe\_maker.sh}).
  \item \code{Unfolder} / \code{UnfolderNuMI} --- extracts and unfolds the
        cross section (\code{run\_unfolder.sh}).
\end{enumerate}

\subsection{Samples}

Run~1 FHC, as configured in \code{configs/file\_properties\_numi.txt}
(Table~\ref{tab:samples}).

\begin{table}[h]
\centering
\small
\begin{tabularx}{\linewidth}{@{}lX@{}}
\toprule
\textbf{Sample type} & \textbf{Purpose} \\
\midrule
\code{onBNB}   & beam-on data (see warning below) \\
\code{extBNB}  & beam-off (EXT), 3\,821\,593 triggers \\
\code{numuMC}  & NuMI overlay central value \\
\code{dirtMC}  & dirt overlay \\
\code{detVarCV} + 8 variations & detector systematics \\
\bottomrule
\end{tabularx}
\caption{Run 1 FHC sample set.}
\label{tab:samples}
\end{table}

\begin{callout}{The \code{onBNB} slot currently holds NuWro fake data}
The active entry is
\code{xsec-ana-numi\_nuwro\_overlay\_pion\_ntuples\_run1\_fhc.root}
($6.65\times10^{20}$ POT). The real beam-on sample
(\code{xsec-ana-neutrinoselection\_filt\_run1\_beamon\_beamgood.root},
$3.283\times10^{20}$ POT, 7\,809\,962 triggers) is \textbf{commented out} at
\code{configs/file\_properties\_numi.txt:31}.

The analysis is therefore configured as a \emph{fake-data study}, not a
real-data measurement. Section~\ref{sec:results} cannot be filled until this
is swapped back and the chain re-run.
\end{callout}

\subsection{Detector variation samples}
\label{sec:detvar}

Eight variations are used: \code{detVarLYdown}, \code{detVarLYrayl},
\code{detVarRecomb2}, \code{detVarSCE}, \code{detVarWMAngleXZ},
\code{detVarWMAngleYZ}, \code{detVarWMX} and \code{detVarWMYZ}.

The BNB note additionally lists \code{detVarLYatten}; it is not present in the
NuMI sample set here. As in the BNB analysis, the WireMod $dE/dx$ variation is
not used, in favour of the recombination variation.

\subsection{Corrections applied since the previous iteration}
\label{sec:corrections}

The following were found and fixed during a code review of this analysis. All
of them affect the numbers a run would produce, which is why no results
predating them should be carried forward.

\begin{itemize}[itemsep=0.3em,topsep=0.3em]
  \item \textbf{Muon momentum estimator.} The estimator combining range
        (contained) and MCS (uncontained) was computed but never written to
        the output tree, so the unfolding binned on MCS alone. It is now
        branched as \code{candidate\_muon\_mom\_reco} and
        \code{ccpi\_pmu\_bin\_config.txt} uses it.
  \item \textbf{Pion momentum estimator.} Reconstructed pion momentum was a
        proton-hypothesis range \emph{kinetic energy} binned against a true
        \emph{momentum} on identical bin edges. It is now the muon-hypothesis
        range momentum.
  \item \textbf{Per-category systematics.} Six \code{MCFullCorrCategory}
        entries evaluated to identically zero and have been removed; see
        Section~\ref{sec:gaps}.
  \item \textbf{Background true bins.} \code{MakeConfig} had stopped emitting
        them, so any regenerated bin configuration would have had no
        background subtraction at all.
  \item \textbf{Diagnostic histograms} are now filled with the central-value
        weight (\code{spline\_weight\_}\,$\times$\,\code{tuned\_cv\_weight\_}).
\end{itemize}

% =============================================================================
\section{Signal and Selection}
\label{sec:selection}

\subsection{Signal definition}

An event is signal if, in truth:

\begin{itemize}[itemsep=0.15em,topsep=0.3em]
  \item the neutrino vertex lies inside the fiducial volume,
        $10 < x < 246$~cm, $-101 < y < 101$~cm, $10 < z < 986$~cm;
  \item the interaction is charged-current;
  \item the neutrino is \numu\ (\code{|mc\_nu\_pdg| == 14}, see
        Section~\ref{sec:open});
  \item there is exactly one muon above threshold;
  \item there is exactly one charged pion;
  \item there are no neutral pions, no kaons and no heavier mesons;
  \item any number of protons is allowed, with no momentum threshold --- the
        ``Xp'' in \code{CC1mu1piXp};
\end{itemize}

and the event falls within the measured phase space:
\begin{equation}
\pmu > 0.15~\GeVc, \qquad
\ppi > 0.175~\GeVc, \qquad
\theta_{\mu\pi} < 2.6~\mathrm{rad}.
\end{equation}

The phase-space cuts are applied identically in \code{define\_signal()} and
\code{categorize\_event()}, so signal and category assignment cannot disagree.

\subsection{Binning}

Five differential observables are measured; bin edges as configured are given
in Table~\ref{tab:binning}.

\begin{table}[h]
\centering
\small
\begin{tabularx}{\linewidth}{@{}llcX@{}}
\toprule
\textbf{Observable} & \textbf{Config} & \textbf{true/reco} & \textbf{Edges} \\
\midrule
\pmu & \code{ccpi\_pmu} & 23/22 &
  0.15, 0.20, 0.30, \ldots, 2.00, 2.30, 2.60, 3.00 \GeVc \\
\ppi & \code{ccpi\_ppi} & 6/5 &
  0.175, 0.20, 0.30, 0.40, 0.50, 1.00 \GeVc \\
$\cos\theta_\mu$ & \code{ccpi\_costhmu} & 13/12 &
  $-1.0$, $-0.5$, 0.0, 0.2, 0.4, 0.55, 0.65, 0.75, 0.82, 0.88, 0.93, 0.97, 1.0 \\
$\cos\theta_\pi$ & \code{ccpi\_costhpi} & 13/12 &
  $-1.0$, $-0.7$, $-0.4$, $-0.2$, 0.0, 0.2, 0.35, 0.5, 0.65, 0.78, 0.88, 0.95, 1.0 \\
$\theta_{\mu\pi}$ & \code{ccpi\_thmupi} & 10/9 &
  0.0, \ldots, 2.513, 2.600 rad \\
\bottomrule
\end{tabularx}
\caption{Binning for the five differential observables.}
\label{tab:binning}
\end{table}

Each true-bin count is one greater than the corresponding reco count: the
extra true bin is the background bin, defined by inverting the signal
definition (\code{!CC1mu1piXp\_MC\_Signal}). All configurations are
single-block with no sideband reco bins.

\subsection{Particle identification}

Three BDTs are used, trained separately and applied via \code{TMVA::Reader}
(Table~\ref{tab:bdt}).

\begin{table}[h]
\centering
\small
\begin{tabularx}{\linewidth}{@{}llX@{}}
\toprule
\textbf{Reader} & \textbf{Weights} & \textbf{Role} \\
\midrule
\code{tmvaReader}    & \code{dataset\_MIP\_BDT\_no\_len}  & MIP identification \\
\code{tmvaReader\_mu} & \code{dataset\_muon\_BDT}          & muon --- \textbf{booked
   but never evaluated}, see Section~\ref{sec:open} \\
\code{tmvaReader\_pi} & \code{dataset\_pion\_BDT\_no\_len} & pion \\
\bottomrule
\end{tabularx}
\caption{Boosted decision trees used for particle identification. Weight files
live under \code{booster\_decision\_tree/}.}
\label{tab:bdt}
\end{table}

Input variables are the three-plane Bragg likelihood ratios
(\code{trk\_bragg\_p\_v}, \code{trk\_bragg\_mu\_v}, \code{trk\_bragg\_mip\_v}),
the LLR PID score, the track score, the track length, and the SCE-corrected
track end position. The pion reader omits the track length.

\subsection{Event selection cuts}

Applied in order; an event passes only if all are satisfied:

\begin{enumerate}[itemsep=0.1em,topsep=0.3em]
  \item software trigger
  \item reconstructed vertex in the fiducial volume
  \item topological score
  \item muon candidate is track-like
  \item pion candidate is contained
  \item muon candidate not in a detector gap
  \item pion candidate not in a detector gap
  \item shower cut
  \item muon--pion opening angle cut
  \item final track-multiplicity cut
\end{enumerate}

\subsection{Selection efficiency and purity}

\emph{To be filled from a complete \code{ProcessNTuples} pass.}
\code{finalize()} prints the efficiency and purity; these are now
central-value weighted (Section~\ref{sec:corrections}) and therefore differ
from any previously recorded values.

% =============================================================================
\section{Uncertainties}
\label{sec:unc}

The uncertainty budget is configured in
\code{configs/ccpi\_systcalc\_numi.conf}. The total is
\begin{equation}
\mathrm{total} = \mathrm{detVar} + \mathrm{flux} + \mathrm{reint}
 + \mathrm{xsec} + \mathrm{POT} + \mathrm{numTargets}
 + \mathrm{SimStats} + \mathrm{DataStats}.
\end{equation}

\subsection{Flux}
PPFX hadron-production multisims via \code{weight\_ppfx\_all}. A 2\,\%
fully-correlated normalisation uncertainty is applied for POT
(\code{POT MCFullCorr 0.02}).

\subsection{Interaction}
GENIE multisims (\code{weight\_All\_UBGenie}) plus nine unisim knobs
(\code{AxFFCCQEshape}, \code{DecayAngMEC}, \code{NormCCCOH},
\code{NormNCCOH}, \code{RPA\_CCQE}, \code{ThetaDelta2NRad},
\code{Theta\_Delta2Npi}, \code{VecFFCCQEshape}, \code{XSecShape\_CCMEC}) and
the two second-class-current form factors (\code{xsr\_scc\_Fa3\_SCC},
\code{xsr\_scc\_Fv3\_SCC}), matching the BNB treatment.

\subsection{Reinteraction}
\code{weight\_reint\_all} multisims.

\subsection{Detector}
The eight detector variation samples of Section~\ref{sec:detvar}, evaluated as
type \code{DV}.

\subsection{Target}
A 1\,\% normalisation uncertainty on the number of argon nuclei in the
fiducial volume (\code{numTargets MCFullCorr 0.01}), as in the BNB analysis.

\subsection{Known gaps in the uncertainty budget}
\label{sec:gaps}

Two components are \textbf{not} included, and the total above is
correspondingly under-estimated.

\textbf{NuMI beamline geometry.} The beamline variation weights
(\code{weight\_Horn\_2kA}, \code{weight\_Horn1\_x\_3mm},
\code{weight\_Beam\_spot\_*} and others) are absent from the processed
ntuples. \code{instructions\_numi.txt} documents
\code{AddBeamlineGeometryWeights} as the tool that injects them, but no driver
script invokes it, and the required histogram file
(\code{NuMI\_Geometry\_Weights\_Histograms.root}, canonical copy at
\code{/exp/uboone/data/users/pgreen/NuMIFlux/NewFluxFiles/}) is not available
locally. This systematic has no BNB counterpart and must be restored before
the measurement is complete.

\textbf{Background normalisation by category, and dirt normalisation.} Six
\code{MCFullCorrCategory} entries were removed because they evaluated to
identically zero: they matched true bins by exact string equality against
\code{"category == N"} while the bin configurations declare a single
\code{"!CC1mu1piXp\_MC\_Signal"} background bin. Their category indices were
additionally inherited from the NuMI CC1e analysis and did not correspond to
this selection's categories. Removing them changed no number --- verified with
152 histograms bit-identical. Full detail is recorded in
\code{configs/ccpi\_systcalc\_numi.README.md}.

For reference, the BNB uncertainty breakdown comprises \code{EXTstats},
\code{MCstats}, \code{POT}, \code{detVar\_total}, \code{flux},
\code{numTargets}, \code{reint} and \code{xsec\_total} --- the same set
retained here, with no per-category or dirt term.

% =============================================================================
\section{Cross-section Extraction}
\label{sec:extraction}

The flux-integrated differential cross section in truth bin $\alpha$ is
obtained by subtracting the background from the measured event rate, unfolding
to truth space, and dividing by the integrated flux $\Phi$, the number of
argon targets $T$, and the bin width:
\begin{equation}
\left.\frac{d\sigma}{dx}\right|_{\alpha}
 = \frac{1}{T\,\Phi\,\Delta x_\alpha}
   \sum_a U_{\alpha a}\left(D_a - B_a\right),
\label{eq:xsec}
\end{equation}
where $D_a$ is the measured event rate in reco bin $a$, $B_a$ the estimated
background, and $U_{\alpha a}$ the unfolding matrix, which absorbs both the
reco-to-truth mapping and the efficiency.

Targets are counted as \emph{argon nuclei} in the fiducial volume, so the
result is per-nucleus rather than per-nucleon, matching the BNB convention
($T = 1.03\times10^{30}$ nuclei there). For NuMI the fiducial volume
additionally excludes the dead region $675.1 < z < 775.1$~cm.

Results are quoted in units of $10^{-38}~\mathrm{cm}^2/\mathrm{Ar}$,
differential in the observable --- so
$10^{-38}~\mathrm{cm}^2/\mathrm{GeV}/\mathrm{Ar}$ for momenta and
$10^{-38}~\mathrm{cm}^2/\mathrm{rad}/\mathrm{Ar}$ for angles.

\begin{callout}{Unresolved --- see Section~\ref{sec:open}}
The code currently divides by $10^{39}$ in NuMI mode and $10^{38}$ otherwise,
and \code{Unfolder.C} does not divide by the bin width while
\code{UnfolderNuMI.C} does. Both must be settled before any number in
Section~\ref{sec:results} is meaningful.
\end{callout}

\subsection{Unfolding}

Wiener-SVD with the filter enabled and second-derivative regularisation
(\code{Unfold WienerSVD 1 second-deriv}), chosen over D'Agostini after a
variant study: it suppresses the bin-to-bin oscillation arising from sparse
Run~1 statistics. D'Agostini with four iterations remains available as a
cross-check via \code{configs/ccpi\_xsec\_config\_numi\_range\_dagost.txt}.

As in the BNB analysis, the unfolded result must be compared to predictions
smeared by the additional smearing matrix $A_C$.

% =============================================================================
\section{Fake-Data Tests}
\label{sec:fakedata}

\emph{Structure only --- no results.}

\subsection{GENIE closure test}
Treat the GENIE central value as data, add the proportionally scaled beam-off
sample, unfold, and confirm that the GENIE CV truth distribution is recovered.
Driven by \code{configs/ccpi\_xsec\_config\_numi.txt}.

\subsection{NuWro fake-data study}
The same procedure with the NuWro overlay
(\code{configs/ccpi\_xsec\_config\_numi\_nuwro.txt}, driven by
\code{run\_nuwro\_observables.sh} across all five observables). As in the BNB
analysis, only statistical and model uncertainties apply, since the samples
differ only in the generator model.

\emph{Per-observable results to be added: reco-space selection histograms,
reco- and truth-space bin correlations, confusion matrices, unfolded cross
sections and additional smearing matrices, for each of $\cos\theta_\mu$,
$\cos\theta_\pi$, \pmu, \ppi\ and $\theta_{\mu\pi}$.}

% =============================================================================
\section{Results}
\label{sec:results}

\emph{Empty pending a real-data run.} This requires, in order: the real
beam-on sample restored in \code{file\_properties\_numi.txt}; the full chain
re-run from \code{ProcessNTuples}, since the corrected momentum branches do
not exist in the current processed files; and Section~\ref{sec:open}, item~1
resolved.

This section should contain, per observable and for the total, the unfolded
cross section with the full uncertainty breakdown, and generator comparisons.

% =============================================================================
\section{Open Items}
\label{sec:open}

\begin{enumerate}[itemsep=0.4em,topsep=0.3em]
  \item \textbf{Cross-section units and bin-width normalisation.}
        \code{conversion\_factor()} divides by $10^{39}$ in NuMI mode against
        $10^{38}$ for BNB, while the BNB note quotes
        $10^{-38}~\mathrm{cm}^2/\mathrm{Ar}$ --- a factor of ten discrepancy
        in convention. Separately, \code{UnfolderNuMI.C} divides by the bin
        width and \code{Unfolder.C} does not, so the two binaries disagree on
        the same bin. The BNB note is unambiguous that the result is
        differential and per-Ar.

  \item \textbf{Beamline geometry systematic.} See Section~\ref{sec:gaps}.

  \item \textbf{Anti-neutrino contamination.} The signal definition uses
        \code{|mc\_nu\_pdg| == 14}, admitting \numubar\ charged-current events
        into both the signal and the efficiency denominator. All three sibling
        selections in the framework use an exact comparison. If the flux
        normalisation is \numu-only, this biases the result.

  \item \textbf{Unused muon BDT.} \code{tmvaReader\_mu} is constructed and
        booked but never evaluated, and its input variables are never filled,
        suggesting that an intended muon-PID cut was dropped.
\end{enumerate}

\vfill
\begin{center}\small
Prepared as a companion to \code{BNBInternalNoteCC1pi\_V0.91.pdf}.\\
Analysis code state: see git history on branch \code{fix/systematics-warning}.
\end{center}

\end{document}
